\chapter{模拟引擎与信号增强}

从 Wilson 作用量到一个物理数字之间隔着机器：一个采样规范场的算法、
一个改进的费米子离散化，以及让强子信号在部分子物理所需的大动量下
可见的涂抹技术。本章四篇论文是这台机器的承重构件；另有两篇
（G\"usken 等 1989 的 Wuppertal 涂抹与 Sheikholeslami--Wohlert 1985
的 clover 改进项）被库内系综使用得如此普遍，在此一并致意而不单独立条。

% ----------------------------------------------------------------------
\section{混合蒙特卡洛（HMC）}
\papercard{S.\ Duane, A.\ D.\ Kennedy, B.\ J.\ Pendleton, D.\ Roweth}
{Hybrid Monte Carlo}
{Phys.\ Lett.\ B \textbf{195}, 216--222 (1987)}
{无 arXiv（印本前时代）| DOI: 10.1016/0370-2693(87)90011-7}
\whyselected{HMC 是本库所用几乎所有动力学系综（MILC HISQ、CLS、CLQCD
clover 系综、domain wall）背后的引擎。库内的 HMC 自动调参论文、两篇
流式采样论文以及粲介子系综论文都明确引用它。}
\physics{HMC 把分子动力学与 Metropolis 联姻：为场变量附加动量 $p$，
沿虚拟哈密顿量轨迹积分长度 $\tau$，
\begin{equation}
  H = \frac{1}{2}p^2 + S[U], \qquad \dot U = i\!\left[\frac{\partial H}{\partial p}\right]\!U ,
\end{equation}
终点以概率 $\min(1,e^{-\Delta H})$ 接受。由于积分子可逆且保体积，
细致平衡严格成立，而全局提议抑制随机游走行为：去相关随 $\tau$ 增长
而非随步数增长。其代价——反复求解 Dirac 方程——推动了三十年的求解器
革新，最终结晶为本库所用的 QUDA 等库中的多重网格与混合精度算法。}
\impact{本库每项格点 PDF 计算的每个系综都由 HMC 或其直接后代
（Hasenbusch 质量预处理、多时间标度）生成。第八章的 AI 采样路线把
自身定义为叠加在同一 Metropolis 逻辑之上的替代提议层——并以本文为其
基线。}
\cardend

% ----------------------------------------------------------------------
\section{解析链涂抹（stout smearing）}
\papercard{C.\ Morningstar and M.\ Peardon}
{$SU(3)$ 链变量在格点 QCD 中的解析涂抹}
{Phys.\ Rev.\ D \textbf{69}, 054501 (2004)}
{arXiv:hep-lat/0307022 | DOI: 10.1103/PhysRevD.69.054501}
\whyselected{stout 涂抹作为流的离散祖先贯穿本库的梯度流章节，出现在
L\"uscher 本人论文的参考文献中（流语料内约 $\sim$7/99）。它也是涂抹
准 PDF 计算的标准降噪工具。}
\physics{细链被向 staple 力做一步指数更新而替换为光滑链，
\begin{equation}
  U_\mu(n)\;\to\; e^{\,\mathrm{i}\rho\,\Omega_\mu(n)}\,U_\mu(n),\qquad
  \Omega_\mu = \mathrm{antiherm}\Big(\sum_{\pm\nu\neq\mu}\!U_\nu(n)U_\mu(n+\hat\nu)U^\dagger_\nu(n+\hat\mu)\Big),
\end{equation}
迭代数次。更新对链变量解析、按构造规范协变且代价低廉。它与梯度流的
联系是本质性的：权重 $\rho$ 的 stout 涂抹正是第三章流方程
（式~\ref{eq:flow}）取流时间步长 $t=\rho$ 的 Euler 离散化。}
\impact{"stout = 粗粒化的流"这一等价性正是 L\"uscher~2010 论证流之
物理含义的方式；Monahan 与 Orginos 在用连续流涂抹准 PDF 算符时利用了
同一结构（第六章）。在本库内它也是核子结构矩阵元的主力涂抹手段。}
\cardend

% ----------------------------------------------------------------------
\section{蒸馏（distillation）}
\papercard{M.\ Peardon, J.\ Bulava, J.\ Foley, C.\ Morningstar, J.\ Juge,
A.\ J\"uttner, K.\ Orginos, et al.\ (Hadron Spectrum Collaboration)}
{用于格点 QCD 强子物理的一种新型夸克场产生算符构造}
{Phys.\ Rev.\ D \textbf{80}, 054506 (2009)}
{arXiv:0905.2160 | DOI: 10.1103/PhysRevD.80.054506}
\whyselected{本库原文之一，也是库内整个 HadStruc 伪分布纲领（胶子
PDF、胶子螺旋度、高动量蒸馏）的使能技术：至少 4 篇库内论文以其为方法
出处。}
\physics{与其在源或汇处临时应用 Jacobi 涂抹，不如一劳永逸地固定一个
由各时间片三维拉普拉斯算符低模构成投影算符，
\begin{equation}
  \square^{[N]}(t)=\sum_{k=1}^{N}\lvert v_k(t)\rangle\langle v_k(t)\rvert,\qquad
  -\square_3\,v_k=\lambda_k v_k,
\end{equation}
并把每个夸克传播子通过 \emph{perambulator} $\tau$ 写出——它是每个
组态只需计算一次的小对象：
$\langle x|e^{-Ht}|y\rangle = \Phi(x,t)\,\tau(t)\,\Phi^\dagger(y,0)$。
此后任意多个插符可以\emph{事后}组合成变分矩阵，无需新的 Dirac 求逆。
L\"uscher--Wolff 变分法因此变得可以规模化使用。}
\impact{在本库内，Egerer 等人的《高动量下的蒸馏》通过附加动量相位把
该思想推进到核子动量 3 GeV 以上——恰是伪 ITD 胶子计算（Khan 等、
Egerer 等）运作的区域。HadStruc 的胶子 PDF 路线离开这篇论文不可想象。}
\cardend

% ----------------------------------------------------------------------
\section{动量涂抹（momentum smearing）}
\papercard{G.\ S.\ Bali, B.\ Lang, B.\ U.\ Musch, A.\ Sch\"afer}
{格点 QCD 中高动量强子的新型夸克涂抹}
{Phys.\ Rev.\ D \textbf{93}, 094515 (2016)}
{arXiv:1602.05525 | DOI: 10.1103/PhysRevD.93.094515}
\whyselected{本语料库中每个 boosted 强子计算的标准信号增强工具：
LaMET 螺旋度、横率、氘核样双重子与运动学增强插符等论文均引用
（约 $\sim$6/99）。}
\physics{普通空间涂抹优化的是与\emph{静止系}波函数的重叠；高速强子的
波函数在振荡，重叠随 $P^z$ 幂次衰减。动量涂抹给涂抹核乘上针对目标
动量调谐的平面波相位，
\begin{equation}
  \tilde\psi(\mathbf{x}) = \sum_{\mathbf{y}} e^{\,\mathrm{i}\mathbf{P}\cdot(\mathbf{x}-\mathbf{y})}\,
  K(\mathbf{x}-\mathbf{y})\,\psi(\mathbf{y}),
\end{equation}
使涂抹后的算符追踪运动态。$P^z\sim1$--$2$ GeV 三点函数的信噪比改善数
倍——常常就是"能提取 PDF"与"不能提取"的分界。}
\impact{在本库内它与蒸馏结合于 HadStruc 系列；Zhang 等的运动学增强
插符论文（报告 $P^2$ 增强重叠再带来量级级增益）与之互补，共同定义了
高动量矩阵元的当前技艺水准。}
\cardend
